\section*{Abstract}

Policy microsimulation models rely on imputation to attach variables one
survey observes to records from another, and the imputed file is then
reweighted downstream: calibrated to administrative totals, filtered to
subpopulations, stressed under reform scenarios. Common practice fits
imputation models to survey records while ignoring their design weights,
which biases imputed conditionals toward the sample rather than the
population. We present the imputation operator of \populace{},
\policyengine{}'s open-source microdata stack: a regime-gated, sequentially
chained, weighted-bootstrap quantile-regression-forest estimator whose
interface makes an unweighted fit impossible to request by accident. To
evaluate it we formalize a \emph{population-view harness}: one latent
population observed through survey views, with any candidate weighted file
scored by projecting it through each view against that survey's holdout.
Benchmarked against standard survey-imputation methods on four task
families, each design choice earns its place where its failure mode lives.
Fit without weights, the identical estimator is six times worse on
net-worth marginal fit and imputes a 99th percentile 1.9 to 2.3 times the
holdout's --- an inflation that capped joint-geometry metrics miss entirely
and only an uncapped tail axis detects. Regime gating holds dividend
zero-share error to 0.6 percentage points against 4.4 for the same forest
ungated. Chaining is worth 5.8 points of classifier two-sample AUC once
four imputed components must satisfy a balance-sheet identity. The
estimator is available standalone as \populacefit{}, and every number in
the paper regenerates from the accompanying repository.

\vspace{1em}
\noindent\textbf{Keywords:} microsimulation, imputation, survey weights,
quantile regression forests, statistical matching
